\documentclass[11pt,a4paper]{article}
\usepackage[italian]{babel}
\usepackage[margin=2.5cm]{geometry}
\usepackage{parskip}
\usepackage{amsmath, amssymb, amsthm}
\usepackage{booktabs}
\usepackage{array}
\usepackage{xcolor}
\usepackage{needspace}
\usepackage{enumitem}
\usepackage{listings}
\usepackage{tikz}
\usepackage[most]{tcolorbox}
\usepackage[hidelinks]{hyperref}
\usetikzlibrary{decorations.pathreplacing, shapes.geometric, arrows.meta, positioning}

% Niente vedove/orfane: i paragrafi non lasciano righe isolate a fine/inizio pagina
\widowpenalty=10000
\clubpenalty=10000
\linespread{0.98}

% Elenchi più compatti
\setlist{itemsep=2pt, parsep=0pt, topsep=4pt, partopsep=0pt}

% --- Riquadri con bordo nero, senza numero (si spezzano tra due pagine) ---
% Uso: \begin{definizione} ... \end{definizione}
%      \begin{teorema}[Nome] ... \end{teorema}  (il nome è facoltativo)
\tcbset{nota/.style={enhanced jigsaw, breakable, colback=white,
  colframe=black, boxrule=0.5pt, sharp corners}}
\NewTColorBox{definizione}{}{nota,
  before upper={\textbf{Definizione.}\ }}
\NewTColorBox{teorema}{o}{nota,
  before upper={\textbf{Teorema\IfValueT{#1}{ (#1)}.}\ }}
\NewTColorBox{proposizione}{o}{nota,
  before upper={\textbf{Proposizione\IfValueT{#1}{ (#1)}.}\ }}
\NewTColorBox{proprieta}{o}{nota,
  before upper={\textbf{Proprietà\IfValueT{#1}{ (#1)}.}\ }}
\NewTColorBox{assioma}{o}{nota, boxrule=1pt,
  before upper={\textbf{Assioma\IfValueT{#1}{ (#1)}.}\ }}

% --- Ambienti semplici (senza riquadro) ---
% Il titolo sta in cima al blocco, anche se il contenuto inizia con un riquadro o
% con codice; e non resta mai da solo in fondo a una pagina.
% Uso: \begin{esempio}[Titolo facoltativo] ... \end{esempio}
\newcommand{\newrunin}[2]{%
  \NewDocumentEnvironment{#1}{o}{%
    \par\Needspace{4\baselineskip}\addvspace{\medskipamount}\noindent
    \textbf{#2}\IfValueT{##1}{ (##1)}\textbf{.}\ \ignorespaces}%
  {\par\addvspace{\medskipamount}}}
\newrunin{esempio}{Esempio}
\newrunin{esercizio}{Esercizio}
\newrunin{osservazione}{Osservazione}

% V e F nelle tavole di verità
\newcommand{\V}{\textsf{V}}
\newcommand{\F}{\textsf{F}}

% --- Codice C ---
% Uso: \begin{codice} ... \end{codice}      codice C numerato
%      \begin{comando} ... \end{comando}    comandi da terminale
%      \begin{risultato} ... \end{risultato} output di un programma
%      \lstinline|printf("ciao");|           codice dentro al testo
\lstdefinestyle{c}{
  language=C,
  basicstyle=\ttfamily\small,
  keywordstyle=\color{blue}\bfseries,
  commentstyle=\color{gray}\itshape,
  stringstyle=\color{red!70!black},
  numbers=left,
  numberstyle=\tiny\color{gray},
  numbersep=8pt,
  columns=flexible,
  keepspaces=true,
  breaklines=true,
  showstringspaces=false,
  tabsize=4,
  morekeywords={include, define},
  % lettere accentate nei commenti
  literate={à}{{\`a}}1 {è}{{\`e}}1 {é}{{\'e}}1 {ì}{{\`i}}1
           {ò}{{\`o}}1 {ù}{{\`u}}1 {È}{{\`E}}1 {À}{{\`A}}1
}
\lstset{style=c}
\newtcblisting{codice}{listing only, nota, left=6mm,
  listing options={style=c}}
\newtcblisting{comando}{listing only, nota,
  listing options={basicstyle=\ttfamily\small, numbers=none, language={},
    columns=flexible, keepspaces=true, breaklines=true}}
\newtcblisting{risultato}{listing only, enhanced jigsaw, breakable,
  colback=black!5, colframe=black, boxrule=0.5pt, sharp corners,
  listing options={basicstyle=\ttfamily\small, numbers=none, language={},
    columns=flexible, keepspaces=true, breaklines=true}}

% --- Diagrammi di flusso (simboli standard) ---
\tikzset{
  terminale/.style = {ellipse, draw, minimum width=2.5cm, minimum height=0.9cm},
  io/.style        = {trapezium, trapezium left angle=70, trapezium right angle=110,
                      draw, minimum width=2.5cm, minimum height=0.9cm},
  processo/.style  = {rectangle, draw, minimum width=2.5cm, minimum height=0.9cm},
  decisione/.style = {diamond, draw, aspect=2, minimum width=2.5cm},
  freccia/.style   = {thick, -{Stealth}}
}

\title{Massimi, minimi, maggioranti e minoranti}
\author{Marco Cerbella}
\date{Lezione 2 -- 28 settembre 2026}

\begin{document}
\maketitle

\section{Massimi e minimi}

\begin{definizione}
Sia $A \subseteq \mathbb{R}$. Un numero $M$ si dice \emph{massimo} di $A$,
e si scrive $M = \max A$, se
\[
M \in A \quad \text{e} \quad a \leq M \ \text{ per ogni } a \in A.
\]
Analogamente, $m$ si dice \emph{minimo} di $A$, e si scrive $m = \min A$, se
\[
m \in A \quad \text{e} \quad m \leq a \ \text{ per ogni } a \in A.
\]
\end{definizione}

Il massimo è quindi l'elemento più grande dell'insieme, il minimo il più
piccolo. La condizione $M \in A$ è essenziale: il massimo deve
\textbf{appartenere} all'insieme.

Se esiste, il massimo è \textbf{unico}: se $M$ e $M'$ fossero entrambi
massimi di $A$, si avrebbe $M \leq M'$ (perché $M'$ è massimo e $M \in A$) e
$M' \leq M$ (per lo stesso motivo), quindi $M = M'$. Lo stesso vale per il
minimo.

\begin{esempio}
\leavevmode
\begin{enumerate}
    \item $A = \{1, 3, 7\}$: $\max A = 7$, $\min A = 1$. Ogni insieme
          \textbf{finito} e non vuoto ha sempre massimo e minimo.
    \item $A = [0, 1] = \{x \in \mathbb{R} \mid 0 \leq x \leq 1\}$:
          $\min A = 0$, $\max A = 1$.
\end{enumerate}
\end{esempio}

\textbf{Non tutti gli insiemi di numeri reali hanno massimo o minimo.}

\begin{esempio}
\leavevmode
\begin{enumerate}
    \item $\mathbb{N}$ ha minimo $0$ ma non ha massimo: per ogni $n$ c'è
          $n + 1$, più grande.
    \item $\mathbb{Z}$ non ha né massimo né minimo.
    \item $A = [0, 1) = \{x \in \mathbb{R} \mid 0 \leq x < 1\}$ ha minimo
          $0$ ma \textbf{non ha massimo}. Il candidato sarebbe $1$, che però
          non appartiene ad $A$. E nessun $x \in A$ è il massimo, perché
          $\frac{x + 1}{2}$ sta ancora in $A$ ed è più grande di $x$.
\end{enumerate}
\end{esempio}


\begin{esempio}[Cinque insiemi di $\mathbb{R}$]
Consideriamo
\[
A = [1, 2], \quad B = \,]3, 7], \quad C = [1, 8] \cup [10, 25[, \quad
D = [1, 2[ \,\cup \{7\}, \quad E = \,]0, 5[.
\]
Sono tutti \textbf{limitati}. Verifichiamo caso per caso.
\begin{itemize}
    \item $A$: $1 \in A$ e $1 \leq x$ per ogni $x \in A$, quindi $1 = \min A$.
          $2 \in A$ e $2 \geq x$ per ogni $x \in A$, quindi $2 = \max A$.
          $A$ ha sia minimo sia massimo.
    \item $B$: $7 \in B$ e $7 \geq x$ per ogni $x \in B$, quindi $7 = \max B$.
          Invece $3 \notin B$, quindi $3$ \textbf{non} è il minimo (e nessun altro
          numero lo è): $B$ ha massimo ma non ha minimo.
    \item $C$: $1 = \min C$, ma $C$ non ha massimo (il candidato $25$ non
          appartiene a $C$): ha minimo ma non massimo.
    \item $D$: $7 \in D$ e $7 \geq x$ per ogni $x \in D$, quindi $7 = \max D$;
          $1 \in D$ e $1 \leq x$ per ogni $x \in D$, quindi $1 = \min D$.
    \item $E$: né minimo né massimo ($0$ e $5$ non appartengono a $E$).
\end{itemize}
\end{esempio}

Il terzo esempio dei casi precedenti mostra che un insieme può ``avvicinarsi'' a un valore senza
raggiungerlo. Per descrivere questa situazione servono i concetti di
maggiorante, minorante ed estremo superiore/inferiore.

\subsection{Maggioranti e minoranti}

\begin{definizione}
Sia $A \subseteq \mathbb{R}$. Un numero $k \in \mathbb{R}$ si dice
\emph{maggiorante} di $A$ se
\[
a \leq k \quad \text{per ogni } a \in A,
\]
e si dice \emph{minorante} di $A$ se
\[
k \leq a \quad \text{per ogni } a \in A.
\]
\end{definizione}

Differenza con massimo e minimo: un maggiorante \textbf{non deve
appartenere} ad $A$. Inoltre, se $k$ è un maggiorante, lo è anche ogni
numero più grande di $k$: i maggioranti, se esistono, sono infiniti.

\begin{esempio}
Per $A = [0, 1)$:
\begin{itemize}
    \item i maggioranti sono tutti i $k \geq 1$, cioè l'insieme $[1, +\infty)$;
    \item i minoranti sono tutti i $k \leq 0$, cioè l'insieme $(-\infty, 0]$.
\end{itemize}
Né $1$ né $2$ appartengono ad $A$, ma sono entrambi maggioranti.
\end{esempio}


\begin{esempio}
Sia $A = [1, 2]$.
\begin{itemize}
    \item $2$ è un maggiorante ($2 \geq x$ per ogni $x \in A$); anche $27$ e
          $15\,000$ sono maggioranti;
    \item $1$ è un minorante ($1 \leq x$ per ogni $x \in A$); anche $-16$ e $0$
          sono minoranti.
\end{itemize}
In generale i maggioranti e i minoranti di un insieme o non ci sono o sono
infiniti.
\end{esempio}

\begin{definizione}
Un insieme $A \subseteq \mathbb{R}$ si dice:
\begin{itemize}
    \item \emph{limitato superiormente} se ammette almeno un maggiorante;
    \item \emph{limitato inferiormente} se ammette almeno un minorante;
    \item \emph{limitato} se è limitato sia superiormente sia inferiormente.
\end{itemize}
\end{definizione}

\begin{esempio}
\leavevmode
\begin{enumerate}
    \item $[0, 1)$ è limitato.
    \item $\mathbb{N}$ è limitato inferiormente (ad esempio da $0$) ma non
          superiormente.
    \item $\mathbb{Z}$ non è limitato né inferiormente né superiormente.
\end{enumerate}
\end{esempio}

Se $A$ ha massimo, allora $\max A$ è un maggiorante di $A$ (ed è l'unico
maggiorante che appartiene ad $A$). Quindi un insieme che ha massimo è
limitato superiormente. Il viceversa è falso: $[0, 1)$ è limitato
superiormente ma non ha massimo.

\end{document}
